\songtitle{Anama'inreesü tüa}

\tag*{2026}\tag{fragment}\tag{wayuu-lyrics}
\begin{notes}
Extract of the intro from \emph{Punta Espada}, from the \emph{Thomas' Entourage} narrative universe.
\end{notes}
\begin{style}
Electroswing with a high-tempo dance beat at 124 BPM in F minor, The track features a prominent brass section with trumpets and trombones playing syncopated staccato stabs, A synthesized upright bass provides a walking bassline that interacts with a four-on-the-floor kick drum and a crisp, high-frequency snare, A bright, processed female vocal performs repetitive melodic hooks with heavy compression and a slight slapback delay, Glitchy record scratches and vinyl crackle effects are layered over the percussion, A swing-style piano plays rhythmic chords on the off-beats, The arrangement uses sudden filter sweeps and momentary silences for rhythmic emphasis, Wayuunaiki, wayuu language
\end{style}
\begin{lyrics}
\songpart{Intro}
\annotation{four-on-the-floor kick drum, vinyl crackle, staccato brass stabs}\\
Anama'inreesü tüa.
\annotation{record scratch}\\
Anama'inreesü tüa.
\annotation{full band enters, walking synth bass, swing piano}\\
Shiimainjeechi ma'in nümaka.

\songpart{Chorus}
\annotation{processed female vocals, bright brass section}\\
Eekaja'a tü sa'awaajünaka ma'in paala; wa'raajüin.

\songpart{Outro}
\annotation{filter sweep on brass}\\
¿Kasa ma'in na'yataajaka anainpünaa waya?
\end{lyrics}

